% 内积空间：由内积定义的距离（三角不等式、诱导范数、平行四边形恒等式、Hilbert空间）。

\section{由内积定义的距离}
\label{sec:ls-inner-product-distance}

\paragraph{三角不等式}
有了长度，自然想用$\lVert\bm x-\bm y\rVert$度量两个向量的距离。这要求长度是一个范数，其中关键是三角不等式：从$\bm x$到$\bm z$的距离不超过经过$\bm y$绕行的距离。定义~\ref{def:ls-normed-space}的正定性直接来自内积的正定性，齐次性来自$\langle c\bm x,c\bm x\rangle=c^2\langle\bm x,\bm x\rangle$。需要证明的是三角不等式，而Cauchy--Schwarz不等式正好控制展开式中的交叉项。
\begin{theorem}[内积诱导长度的三角不等式]
\label{thm:ls-inner-product-triangle}
设$V$是内积空间，$\lVert\bm v\rVert=\sqrt{\langle\bm v,\bm v\rangle}$，则对任意$\bm x,\bm y\in V$，
\begin{equation}
 \lVert\bm x+\bm y\rVert\leq\lVert\bm x\rVert+\lVert\bm y\rVert,
 \label{eq:ls-inner-product-triangle}
\end{equation}
等号成立当且仅当其中一个向量是另一个向量的非负实数倍。
\end{theorem}
\begin{proof}
按双线性与对称性展开平方，再对交叉项用$\langle\bm x,\bm y\rangle\leq|\langle\bm x,\bm y\rangle|$与Cauchy--Schwarz不等式~\eqref{eq:prelim-cauchy-schwarz}：
\[
 \begin{aligned}
 \lVert\bm x+\bm y\rVert^2
 &=\lVert\bm x\rVert^2+2\langle\bm x,\bm y\rangle+\lVert\bm y\rVert^2\\
 &\leq\lVert\bm x\rVert^2+2\lVert\bm x\rVert\lVert\bm y\rVert+\lVert\bm y\rVert^2\\
 &=(\lVert\bm x\rVert+\lVert\bm y\rVert)^2.
 \end{aligned}
\]
两边都非负，开平方即得不等式。

等号成立当且仅当$\langle\bm x,\bm y\rangle=\lVert\bm x\rVert\,\lVert\bm y\rVert$。若$\bm y=\bm0$，结论显然。若$\bm y\neq\bm0$，由Cauchy--Schwarz不等式的等号条件，$\bm x=s\bm y$；代入得$s\lVert\bm y\rVert^2=|s|\,\lVert\bm y\rVert^2$，即$s\geq0$。反过来，$\bm x=s\bm y$且$s\geq0$时，两侧都等于$(1+s)\lVert\bm y\rVert$。
\end{proof}

\paragraph{内积诱导的范数与距离}
由定理~\ref{thm:ls-inner-product-triangle}，长度$\lVert\bm v\rVert=\sqrt{\langle\bm v,\bm v\rangle}$满足范数的三条要求，称为内积\emph{诱导的范数}。它给出距离
\begin{equation}
 d(\bm x,\bm y)=\lVert\bm x-\bm y\rVert=\sqrt{\langle\bm x-\bm y,\bm x-\bm y\rangle},
 \label{eq:ls-inner-product-distance}
\end{equation}
欧氏距离与Frobenius距离都是这样得到的。于是每个内积空间都是赋范空间，前三节关于范数、收敛与完备性的结论都适用于它。

与一般范数相比，这个距离还可以用内积展开：
\[
 \lVert\bm x-\bm y\rVert^2=\lVert\bm x\rVert^2-2\langle\bm x,\bm y\rangle+\lVert\bm y\rVert^2,
\]
距离的平方由两个长度与夹角决定，这正是余弦定理。下面的平行四边形恒等式说明，这一额外结构可以用来区分哪些范数来自内积。

\paragraph{平行四边形恒等式}
反过来，并非每个范数都由内积给出。内积诱导的范数还满足一条额外的等式，即\term{平行四边形恒等式}。几何上，它表示平行四边形两条对角线长度的平方和等于四条边长度的平方和。
\begin{theorem}[平行四边形恒等式]
\label{thm:linear-algebra-parallelogram}
设$V$是实内积空间，$\bm x,\bm y\in V$，以该内积诱导的范数计量长度，则
\begin{equation}
 \lVert\bm x+\bm y\rVert^2+\lVert\bm x-\bm y\rVert^2
 =2\lVert\bm x\rVert^2+2\lVert\bm y\rVert^2.
 \label{eq:functional-parallelogram-identity}
\end{equation}
欧氏范数是这一结论的特例。
\end{theorem}
\begin{proof}
记$\lVert\bm u\rVert^2=\langle\bm u,\bm u\rangle$。由内积的双线性和对称性，分别展开和向量与差向量的长度平方：
\[
 \begin{aligned}
 \lVert\bm x+\bm y\rVert^2
 &=\lVert\bm x\rVert^2+2\langle\bm x,\bm y\rangle+\lVert\bm y\rVert^2,\\
 \lVert\bm x-\bm y\rVert^2
 &=\lVert\bm x\rVert^2-2\langle\bm x,\bm y\rangle+\lVert\bm y\rVert^2.
 \end{aligned}
\]
两式相加，交叉项相互抵消，得到所需恒等式。证明只用了内积公理，因此适用于定义~\ref{def:linear-algebra-inner-product}中的任一实内积。
\end{proof}

这条恒等式可以用来检验一个范数是否来自内积：只要某个范数在一对向量上使恒等式两侧不等，它就不可能由任何内积给出。下面用这个办法检验$\ell_1$与$\ell_\infty$范数。
\begin{example}[反例：$\ell_1$与$\ell_\infty$范数不来自内积]
\label{ex:ls-l1-linf-not-inner-product}
取$\mathbb R^2$中的两个标准坐标向量$\bm x=\bm e_1$、$\bm y=\bm e_2$作为平行四边形的两条边，两条对角线分别是$\bm x+\bm y=(1,1)^\top$与$\bm x-\bm y=(1,-1)^\top$。在$\ell_1$、$\ell_2$、$\ell_\infty$范数下，四条边的长度都是$1$，差别出在对角线上。

对$\ell_1$范数，两条对角线都长$2$：
\[
 \begin{gathered}
 \lVert\bm x+\bm y\rVert_1^2+\lVert\bm x-\bm y\rVert_1^2=2^2+2^2=8,\\
 2\lVert\bm x\rVert_1^2+2\lVert\bm y\rVert_1^2=2\cdot1^2+2\cdot1^2=4.
 \end{gathered}
\]
对$\ell_\infty$范数，两条对角线都只长$1$：
\[
 \begin{gathered}
 \lVert\bm x+\bm y\rVert_\infty^2+\lVert\bm x-\bm y\rVert_\infty^2=1^2+1^2=2,\\
 2\lVert\bm x\rVert_\infty^2+2\lVert\bm y\rVert_\infty^2=4.
 \end{gathered}
\]
$\ell_1$范数使左侧偏大，$\ell_\infty$范数使左侧偏小，恒等式都不成立。所以不存在内积使$\lVert\bm x\rVert_1=\sqrt{\langle\bm x,\bm x\rangle}$对一切$\bm x$成立；$\ell_\infty$范数同样不是由内积给出的长度。

作为对照，欧氏范数下两条对角线长都是$\sqrt2$，左侧为$4$，与右侧相等。
\end{example}
因此一般范数只能度量大小，不能用上面的公式定义夹角。

另一方面，内积可以由它诱导的范数恢复。在实空间中，分别展开$\lVert\bm x+\bm y\rVert^2$与$\lVert\bm x-\bm y\rVert^2$并相减，得到极化恒等式
\begin{equation}
 \langle\bm x,\bm y\rangle=\frac14\bigl(\lVert\bm x+\bm y\rVert^2-\lVert\bm x-\bm y\rVert^2\bigr).
 \label{eq:ls-polarization-identity}
\end{equation}
所以两个实内积若给出同一个范数，就必相同。

\paragraph{Hilbert空间}
\label{sec:ls-hilbert-spaces}
内积诱导的范数使内积空间成为赋范空间。若这个范数下的每个柯西列都收敛到空间内的元素，就称它为\term{Hilbert空间}（Hilbert space，也译作希尔伯特空间）。

由定理~\ref{thm:ls-finite-dimensional-complete}，有限维内积空间都是Hilbert空间。后面的正交分解与正交投影都用有限维的代数构造得到，不依赖完备性。无限维空间中，完备性与子空间的闭性都需要另行检查；无限正交级数、闭子空间上的投影以及连续泛函的表示，属于本章开头提到的无限维理论。
